% ==============================================================================
% JUSTIFICACIÓN
% ==============================================================================
% Contenido redactado y provisto directamente por el estudiante (no adaptado
% del protocolo; reestructura las subsecciones 3.6-3.9 del protocolo
% "Utilidad para..." en una organización más estándar de tesis: pertinencia,
% relevancia técnica, relevancia social/económica, relevancia institucional,
% beneficiarios). No se modificó el contenido, solo se revisó.
%
% NOTA: el protocolo incluía una subsección explícita "3.7 Utilidad para el
% estudiante" que este capítulo no reproduce como sección propia. Puede ser
% una omisión involuntaria o una decisión editorial (es una categoría de
% justificación más débil/subjetiva que las demás); no se asume ninguna de
% las dos, se deja señalado.
%
% Verificación de fuentes (2026-08-08): se confirmó por WebFetch que
% soy502fortinet2026 respalda literalmente la cifra de 726 millones de
% intentos de ataque citada (dato de 2025). No se pudo verificar
% mastercard2025 (la URL devuelve 403 Forbidden); la cifra de 200% de
% incremento y el ranking "top 10 en Latinoamérica" quedan sin confirmar por
% este medio y deben revisarse manualmente antes de la entrega final.
% Actualización 2026-10-09: zeng2024 se reemplazó por zhen2025anomaly (mismo
% documento IEEE, autores y año corregidos). Se retiró la cita a verma2020 de la
% afirmación sobre el costo de los SIEM: ese artículo es una revisión de IDS
% para IoT y no respalda la afirmación. Resuelto el mismo día: la oración se
% reformuló sobre enisa2021smes (verificado en el PDF: págs. impresas 16, 17 y
% 18); se retiró la afirmación sobre "Machine Learning integrado en SIEM" e
% "infraestructura de análisis masivo", que no tiene fuente verificada.
% Actualización 2026-10-10: se acotó la promesa de entrega del dataset a scripts y
% vectores derivados (dos fuentes, XSS_dataset.csv y command injection.csv, no
% tienen licencia especificada por su origen). Estado verificado ese día: los
% archivos de vectores (training/data_clean/*.parquet|csv, training/splits/*.parquet)
% están en .gitignore y NO están publicados en el repositorio; ver la
% recomendación correspondiente en n-recomendaciones.tex.
\ifdefined\modoensamble
\else
	\input{0-preambulo_standalone}
\fi

\chapter{Justificación}

\section{Importancia y pertinencia del problema}
Como se estableció en el capítulo de Antecedentes, la carencia identificada no es la ausencia de
investigación sobre detección de ataques web mediante aprendizaje automático, sino la ausencia
de una solución que integre esa detección de forma nativa en el runtime de aplicaciones Node.js,
cubra múltiples vectores de ataque en un único sistema y sea operable por organizaciones sin
personal de seguridad dedicado. Justificar este trabajo exige, entonces, explicar por qué esa
carencia específica importa y para quién.

Las aplicaciones web construidas sobre Node.js constituyen una superficie de ataque
particularmente expuesta: el ecosistema npm se caracteriza por cadenas extensas de dependencias transitivas: un estudio
empírico publicado en 2022 reportó que el 42.1~\% de los paquetes analizados tenía al menos una
dependencia directa o indirecta vulnerable en su última versión \cite{zerouali2022impact}, y otro de
2019 concluyó que paquetes individuales pueden afectar a gran parte del ecosistema y que un número
muy reducido de cuentas de mantenedores podría usarse para inyectar código malicioso en la mayoría
de los paquetes \cite{zimmermann2019small}. Además, la ausencia de controles a nivel de aplicación
deja expuesta la lógica de negocio incluso cuando existen defensas perimetrales. Los mecanismos de defensa tradicionales, sea un Web Application Firewall (WAF)
o un Sistema de Detección de Intrusiones (IDS), fueron diseñados para operar en el borde de la red
y dependen de firmas estáticas; esta dependencia produce una tasa de falsos negativos
estructuralmente alta frente a variantes ligeramente modificadas de un ataque conocido o frente a
abuso de la lógica de negocio de la aplicación, categorías que no son detectables por coincidencia de patrones. Esta limitación es coherente con la
comparación experimental de Sureda Riera et al., en la que las herramientas RASP evaluadas superaron
a los WAF \cite{bermejohiguera2022}. Este no es un problema exclusivamente técnico: es un problema de
acceso. Las organizaciones que sí pueden mitigar esta limitación suelen hacerlo mediante plataformas de
gestión de eventos e información de seguridad (SIEM). Para las pequeñas y medianas empresas, sin
embargo, estos sistemas representan inversiones considerables, y muchas soluciones de
ciberseguridad requieren conocimiento especializado para implementarse y administrarse; en el
estudio de la Agencia de la Unión Europea para la Ciberseguridad sobre pymes europeas, solo el
28~\% de los participantes de la encuesta declaró haber asignado el rol de responsable de seguridad
de la información \cite{enisa2021smes}. Aunque ese estudio no se refiere a Guatemala, describe las
mismas dos barreras, presupuesto y personal especializado, que enfrentan las organizaciones que
constituyen la población objetivo de este trabajo (Sección~\ref{sec:just-contexto}). El problema, por tanto, no se limita a la existencia de ataques contra
aplicaciones Node.js, sino a la distribución desigual de las herramientas capaces de detectarlos.

\section{Relevancia técnica y científica}
Desde el punto de vista técnico, este trabajo es pertinente porque aborda una limitación
metodológica recurrente en la literatura de detección de ataques web mediante aprendizaje
automático: la evaluación de modelos sobre payloads o cadenas de texto aisladas, sin considerar
las restricciones de un entorno de ejecución real. Los sistemas de detección basados
exclusivamente en reglas colapsan frente a cargas útiles polimórficas o sintaxis de ataque
ligeramente ofuscada que evade patrones conocidos, mientras que los enfoques de aprendizaje
automático permiten generalizar el comportamiento base del servidor y detectar desviaciones
estadísticas sin depender de actualizaciones manuales de firmas \cite{zhen2025anomaly}. La combinación
de un componente supervisado, orientado a clasificar con precisión los vectores de ataque
conocidos, con un componente no supervisado orientado a identificar anomalías estadísticas
asociadas a variantes no observadas durante el entrenamiento, se propone responder a esta limitación sin comprometer la capacidad de generalización del sistema frente a amenazas nuevas; en qué medida lo logra se evalúa en los capítulos de Resultados y de Análisis de resultados.

La relevancia científica del proyecto radica en que el problema de integrar inferencia de
aprendizaje automático dentro de un proceso Node.js sin degradar su rendimiento no está resuelto
de forma satisfactoria en la literatura revisada, precisamente porque la mayoría de los sistemas
comparables opera como componente externo (proxy inverso o plataforma SIEM) y no como
extensión del propio proceso de la aplicación. Ejecutar inferencia dentro del ciclo de vida de una
petición HTTP exige resolver un reto de ingeniería específico de las arquitecturas orientadas a
eventos: evitar que el cómputo del modelo bloquee el bucle de eventos (\textit{event loop}) del cual
depende toda la capacidad de manejar peticiones concurrentes de Node.js. La construcción de un
pipeline de extracción de características e inferencia que opere de forma asíncrona sobre datos de
naturaleza multidimensional, sin bloquear dicho bucle, constituye un aporte técnico verificable
independientemente del desempeño final del modelo de clasificación, y es generalizable a
cualquier otro problema de seguridad en tiempo de ejecución sobre arquitecturas orientadas a
eventos.

\section{Relevancia social y económica en el contexto guatemalteco}
\label{sec:just-contexto}
La pertinencia de este trabajo no puede argumentarse en abstracto: depende directamente del
contexto en que se desarrollaría la solución. En Guatemala, los ataques cibernéticos registrados
durante 2024 aumentaron un 200~\% respecto del año anterior, cifra que posicionó al país entre los
diez de mayor exposición cibernética en Latinoamérica según el Cyber Insights Report de
Mastercard \cite{mastercard2025}. Reportes posteriores muestran que los incidentes continuaron: durante 2026 se reportaron ataques
cibernéticos contra instituciones públicas guatemaltecas, entre ellos el ataque a la Dirección
General de Control de Armas y Municiones (Digecam) ocurrido el 7 de abril \cite{emisoras2026digecam}
y el ataque cibernético confirmado por la Universidad de San Carlos de Guatemala
\cite{emisoras2026usac}, en un país que, además, no cuenta todavía con un marco legal de
ciberseguridad vigente que module la respuesta institucional a estos incidentes
\cite{lahora2026estancada}.

Este panorama de riesgo creciente contrasta con la estructura del tejido empresarial
guatemalteco. Según cifras del Ministerio de Economía, las micro, pequeñas y medianas empresas
representaron el 99.57~\% de las empresas activas con ventas en el año base 2015
\cite{mineco2015mipyme}, y una nota de la Agencia Guatemalteca de Noticias de julio de 2025 afirma
que representan el 96.3~\% de las empresas guatemaltecas \cite{agn2025mipymes}, y
es razonable asumir que una proporción sustancial de las empresas de base tecnológica que
desarrollan servicios sobre Node.js pertenece a este mismo segmento, es decir, opera sin un
Centro de Operaciones de Seguridad ni presupuesto destinado a licencias de plataformas
empresariales de detección. Nótese que esta es una inferencia razonable a partir de la estructura
general del tejido empresarial del país, y no una cifra específica y verificada sobre el segmento de
startups tecnológicas, dato que no está disponible públicamente de forma desagregada; el
argumento de relevancia social se sostiene sobre la composición general del sector empresarial
guatemalteco, no sobre una estadística exclusiva del sector de software que actualmente no existe.
En ese contexto, una herramienta que se distribuye sin costo mediante un gestor de paquetes de
uso estándar y que no exige modificar la infraestructura de red reduce directamente las dos
barreras de adopción más frecuentemente documentadas para estas organizaciones: el presupuesto y la complejidad operativa \cite{enisa2021smes}.

\section{Relevancia institucional}
Para la Universidad del Valle de Guatemala, este trabajo se ubica en la intersección entre Ciencia
de Datos, Ciberseguridad e Inteligencia Artificial, tres áreas de demanda creciente en el mercado
tecnológico regional. La metodología documentada de construcción del corpus (auditoría de fuentes, etiquetado
determinista y partición verificada por hash) y los vectores de características derivados
constituyen un recurso académico que puede emplearse en cursos de seguridad informática y
aprendizaje automático dentro de la Facultad de Ingeniería,
independientemente de si el sistema resultante alcanza los criterios de aceptación definidos para
el proyecto. La alineación explícita del diseño con marcos de referencia reconocidos
internacionalmente, como OWASP Top 10 y MITRE ATT\&CK, permite además que la metodología
documentada sea replicada, contrastada o extendida por trabajos de graduación posteriores dentro
del mismo programa.

\section{Beneficiarios directos e indirectos}
Es posible distinguir con claridad entre los beneficiarios directos de este trabajo y sus beneficiarios
indirectos, distinción necesaria para no sobreestimar el alcance real de la investigación en esta
etapa.

Los beneficiarios directos son las organizaciones guatemaltecas que desarrollan aplicaciones
sobre Node.js sin equipo de seguridad dedicado y que podrían instalar la librería resultante como
un componente adicional de su arquitectura existente sin incurrir en costo de licenciamiento ni en
modificaciones estructurales de su infraestructura de red. Dentro de este grupo se incluyen
también iniciativas comunitarias, organizaciones no gubernamentales con presencia digital y
proyectos de gobierno abierto que comparten la misma restricción de recursos técnicos y
financieros que las startups tecnológicas.

Los beneficiarios indirectos son de tres tipos. Primero, la comunidad académica de la
Universidad del Valle de Guatemala, a través de la metodología documentada de construcción del corpus, de los vectores derivados y del precedente
metodológico que este trabajo establece para investigaciones futuras en la intersección de
ciberseguridad e inteligencia artificial aplicada. Segundo, otros investigadores y equipos de desarrollo de la región centroamericana que, a partir del
material publicado del proyecto, podrían reproducir la construcción del corpus desde las mismas
fuentes públicas y reutilizar los vectores de características derivados para entrenar o validar
sistemas de detección propios, en un contexto donde la investigación aplicada en ciberseguridad con
datos representativos de la región es escasa. El proyecto no redistribuye los datos crudos de las
fuentes cuya licencia no está especificada por su origen: pone a disposición los analizadores y
scripts de normalización, el archivo de mapeo de etiquetas, la auditoría de fuentes y los
vectores de características derivados, que contienen valores numéricos y no el texto de las
peticiones. La reproducción del corpus no es completa: los conjuntos sintéticos de Command
Injection, que aportan el 17.8~\% de las filas de esa clase en el conjunto de prueba, no cuentan
con un generador versionado. Tercero, y de forma
más amplia, el ecosistema tecnológico guatemalteco en su conjunto, en la medida en que la
adopción voluntaria de controles de seguridad accesibles contribuye a elevar el nivel general de
madurez en ciberseguridad del país mientras el marco regulatorio correspondiente continúa en
trámite legislativo \cite{lahora2026estancada}.

Es importante precisar que la magnitud real de estos beneficios, tanto en términos de tasa de
detección como de impacto en el rendimiento del servidor anfitrión, es precisamente lo que este
trabajo se propone medir y no puede darse por establecida antes de la fase de evaluación descrita
en la metodología. Esta justificación argumenta la pertinencia de investigar el problema, no
anticipa el desempeño del sistema que se construirá para abordarlo.

\ifdefined\modoensamble
\else
	\section*{Referencias}
	\printbibliography[heading=none]
	\end{document}
\fi
